\subsection{Revenue and expense recognition}

These entries cover when revenue and expenses are recognised and how capitalisation shifts expense between periods and cash-flow sections.

\begin{qfmethod}[Revenue recognition: the contract-based model of IFRS~15 and Topic~606, timing and measurement choices, contract balances, gross versus net reporting, and implications for analysis]\label{meth:revenue-recognition-analysis}
\qffield{Statement} \emph{Core principle.} Revenue depicts the transfer of promised goods or services to customers at the amount of consideration to which the entity expects to be entitled \citep[IFRS~15 para.~2]{eu-2023-1803-ifrs15}. The model has five steps:
\begin{enumerate}
\item identify a contract: approved, with identifiable rights and payment terms, commercial substance, and collection of the consideration probable (para.~9);
\item identify each distinct good or service as a performance obligation (para.~22);
\item measure the transaction price, excluding amounts collected for third parties (para.~47); variable consideration (discounts, rebates, refunds, bonuses, penalties) is estimated (paras.~50--51) and included only to the extent that it is highly probable that a significant reversal will not occur (para.~56);
\item allocate the price to performance obligations by relative stand-alone selling price (paras.~73--74);
\item recognise revenue when or as control passes (para.~31): over time if the customer consumes the benefit as the entity performs, controls the asset as it is created, or the asset has no alternative use and the entity has an enforceable right to payment for performance to date (para.~35); otherwise at a point in time, judged by indicators such as a present right to payment, legal title, physical possession, risks and rewards, and acceptance (para.~38).
\end{enumerate}

\emph{Balances and disclosures.} A contract liability arises when the customer pays before performance and a contract asset when performance precedes an unconditional right to payment; an unconditional right is a receivable (paras.~105--108). Incremental costs of obtaining a contract that are expected to be recovered are capitalised (para.~91). Entities disaggregate revenue by factors affecting its nature, timing and uncertainty, and disclose the price allocated to unsatisfied obligations with its expected timing (paras.~114 and~120). An entity that controls a good before transfer is a \emph{principal} and reports gross revenue; an agent reports only its fee or commission (paras.~B35--B37).

\emph{US GAAP.} Topic~606 was developed jointly with IFRS~15 and reaches the same conclusions on all requirements with minor differences. The word ``probable'' in the collectability threshold has different meanings in the two frameworks. Topic~606 prohibits reversing impairment of capitalised contract costs, which IFRS~15 requires. Interim disclosure requirements also differ \citep[pp.~10--11]{us-revenue-2014}. Public business entities apply it to annual periods beginning after 15~December 2017, after a one-year deferral \citep[p.~1]{us-revenue-deferral-2015}; IFRS~15 applies from 1~January 2018 \citep[IFRS~15 para.~C1]{eu-2023-1803-ifrs15}.

\qffield{Notation} SSP stand-alone selling price; TP transaction price; allocation to obligation $j$: $\mathrm{TP}\times \mathrm{SSP}_j/\sum_i\mathrm{SSP}_i$.

\qffield{Conditions} IFRS~15 as adopted in the EU in 2023. The constraint's ``highly probable'' has no numerical threshold in the standard; the 90~percent below is an illustration. Licences, contract modifications and repurchase agreements were not reviewed here.

\qffield{Practice} \emph{Original.} A customer pays 1{,}000 on 1~October 2025 for equipment (SSP 900), installation (SSP 100) and 24 months of service (SSP 200).
\begin{itemize}
\item Allocation gives 750.00, 83.33 and 166.67.
\item The equipment is delivered on 1~November and installed by 30~November; service runs from 1~December. Revenue is 840.28 in 2025, 83.33 in 2026 and 76.39 in 2027, together the full price.
\item The contract liability at 31~December 2025 is 159.72.
\item A volume rebate makes the unit price 100, 95 or 90 with probabilities 0.5, 0.3 and 0.2. The expected value is 96.50, but only 90 is at least 90~percent likely not to reverse, so 90 enters the transaction price until the uncertainty resolves.
\item A marketplace selling a third party's goods for 100 and paying the vendor 85 reports revenue of 100 as principal or 15 as agent. Profit is 15 either way, but the margin is 15 against 100~percent, and revenue-based multiples and turnover change.
\end{itemize}

\emph{Analysis implications.} Revenue judgements (timing of control, estimates of returns and rebates, principal or agent) are where earnings management concentrates \citep[pp.~110--111]{financial-statement-analysis-2009}. Check them in this order:
\begin{itemize}
\item growth in receivables and contract assets against revenue;
\item movements in contract liabilities, which give advance billings visibility;
\item the backlog of unsatisfied obligations;
\item changes in disaggregated revenue mix.
\end{itemize}

\qffield{Pitfalls} Comparing gross-reporting and net-reporting businesses on revenue multiples or margins. Reading a fall in contract liabilities as weaker demand without checking billing terms. Treating bill-and-hold or channel-loading shipments as sales without evidence that control passed (\hyperref[meth:reporting-warning-signs]{warning signs}).

\qffield{See also} \hyperref[meth:expense-recognition-capitalisation]{expense recognition}; \hyperref[meth:earnings-management-presentation]{earnings management}.

\end{qfmethod}

\begin{qfmethod}[Expense recognition and capitalisation: matching through inventories and depreciation, research and development, borrowing costs and software, and how capitalising rather than expensing changes profit, cash-flow sections, assets and ratios]\label{meth:expense-recognition-capitalisation}
\qffield{Statement} \emph{Matching through assets.} When inventories are sold their carrying amount is expensed in the period the related revenue is recognised \citep[IAS~2 para.~34]{eu-2023-1803-ias2}. An item of property, plant and equipment is capitalised only if future benefits are probable and its cost is reliably measurable \citep[IAS~16 para.~7]{eu-2023-1803-ias16}. Day-to-day servicing is expensed, while replacement parts that meet the recognition principle are capitalised (paras.~12--13). The depreciable amount is allocated systematically over the useful life; lives, residual values and methods are reviewed each year and revised as changes in estimate (paras.~50--51 and~61). Borrowing costs directly attributable to a qualifying asset are capitalised \citep[IAS~23 para.~8]{eu-2023-1803-ias23}.

\emph{Research and development.} Under IFRS, research is expensed \citep[IAS~38 para.~54]{eu-2023-1803-ias38}. Development is capitalised only when the entity can demonstrate six conditions: technical feasibility, intention, ability to use or sell, probable economic benefits, adequate resources, and reliable measurement (para.~57). Expenditure first expensed is never later capitalised (para.~71). US GAAP expenses all research and development costs within its scope \citep[para.~12]{us-research-development-1974}. Software to be sold is capitalised only after technological feasibility is established \citep[paras.~3--4]{us-software-1985}. For internal-use software, from annual periods beginning after 15~December 2027 capitalisation starts once management has authorised and committed funding and completion is probable \citep[pp.~2--3]{us-software-2025}.

\emph{Capitalise or expense.} The total cost and total cash are the same either way. Capitalising changes four things:
\begin{itemize}
\item timing: profit is higher while spending exceeds amortisation;
\item cash-flow section: the outlay becomes investing rather than operating;
\item asset base: assets and equity are larger;
\item ratios: interest coverage rises when borrowing costs are capitalised.
\end{itemize}
Delaying expense recognition by capitalising costs is a classic income-shifting device \citep[pp.~110--111]{financial-statement-analysis-2009}.

\qffield{Notation} With annual spend $x_t$ amortised straight-line over $L$ years from the year after spending, the pre-tax profit difference is $x_t-\sum_{k=1}^{L}x_{t-k}/L$. The cumulative difference equals the capitalised asset balance.

\qffield{Conditions} The comparison holds total cash constant and ignores tax unless stated; with a tax deduction for the spending, capitalisation also creates a deferred tax liability. US GAAP references are to the issuing Statements and Updates, not the current Codification.

\qffield{Practice} \emph{Original.} Development spending of 30 a year is amortised over three years.
\begin{itemize}
\item With constant spending, the profit uplift from capitalising is 30, 20 and 10 in years 1--3, then zero once amortisation catches up. The asset settles at 60.
\item With spending growing 10~percent a year, the uplift persists: 8.26 in year~6, with an asset of 90.91. The cumulative uplift equals that asset.
\item Operating cash flow is 30 higher and investing cash flow 30 lower in every year. Total cash is unchanged.
\item Capitalising 8 of 20 borrowing costs on an asset under construction lifts reported interest cover from 3.0 to 5.0 times. Cash interest is unchanged.
\end{itemize}
Harbour capitalised development of 22 and amortised 12 in 2025 (\hyperref[ex:harbour-tools-statements]{case}). Expensing it would cut 2025 operating profit by 10 and operating cash flow by 22.

\qffield{Pitfalls} Comparing operating cash flow or EBITDA between a capitaliser and an expenser without moving the capitalised outlay to the same section. Reading a rising capitalised balance as investment success without checking amortisation lives and impairments. Applying the IAS~38 conditions to US GAAP statements.

\qffield{See also} \hyperref[meth:intangible-acquisition-routes]{intangible assets}; \hyperref[meth:accounting-analysis-adjustments]{comparability adjustments}.

\end{qfmethod}
